\chapter{Wyniki i analiza}
\label{ch:wyniki}

\section{Poprawność i wydajność wariantów}
\label{sec:poprawnosc-wydajnosc}

Główny wariant programu zarządza buforami operatorem zamiany \texttt{<=>}.
Rozważono też wariant naprzemiennego buforowania (ping-pong). Ponieważ oba nie
różnią się znacząco wydajnością, dalsze wyniki dotyczą wariantu głównego
z~operatorem zamiany.

\subsection{Weryfikacja poprawności numerycznej}
\label{subsec:poprawnosc}

Przed przystąpieniem do jakichkolwiek pomiarów wydajnościowych konieczne było
ustanowienie wiarygodnego kryterium poprawności obliczeń. Bez takiego kryterium
przyspieszenie z~równoległości nie ma znaczenia, łatwo
bowiem uzyskać krótszy czas wykonania kosztem cichego naruszenia poprawności
wyniku. W~przypadku badanego programu przyjęto niezmiennik oparty na
własnościach pola temperatury oraz na jego całkowitej sumie.

Warunki początkowe i~brzegowe zadania dobrano tak, aby rozwiązanie pozostawało
przez cały czas symulacji ograniczone do przedziału $[1,2]$. Program po
zakończeniu obliczeń raportuje statystyki końcowego pola w~postaci wpisu
\texttt{final field: min=1.0 max=2.0 sum=<S>}. Utrzymanie dolnej granicy
\num{1.0} oraz górnej granicy \num{2.0} jest bezpośrednią konsekwencją zasady
maksimum dla równania przewodnictwa cieplnego: schemat jawny z~odpowiednio
dobranym krokiem czasowym nie może wygenerować wartości spoza zakresu
wyznaczonego przez dane początkowe i~brzegowe. Naruszenie któregokolwiek
z~tych ograniczeń natychmiast sygnalizowałoby błąd implementacji schematu
różnicowego lub niestabilność numeryczną.

Kluczową rolę w~weryfikacji odgrywa jednak suma wszystkich elementów pola,
w~przybliżeniu równa liczbie komórek siatki. Jednocześnie suma $S$ dobrze
oddaje stan całej dziedziny i~szybko reaguje na~różnice między konfiguracjami:
najmniejsze rozejście się wyników widać jako rozbieżność sum.

Właśnie na tej podstawie zweryfikowano poprawność obliczeń wielowęzłowych.
Przebieg uruchomiony na pojedynczym locale (\texttt{-nl 1}) daje pewną wartość
sumy odniesienia. Ten sam problem rozdzielony na $N$~locale daje identyczną
sumę \emph{pod warunkiem}, że po każdym kroku czasowym wykonywana jest
wymiana warstw brzegowych (halo) za pomocą operacji \texttt{updateFluff}.
Pominięcie tej wymiany skraca czas wykonania, lecz suma $S$ zaczyna się
rozchodzić względem wartości odniesienia. Wynik staje się \emph{cicho błędny}.
W~przeprowadzonych pomiarach (dziedzina $1000^3$, 100 kroków) suma końcowa pola
wynosiła $1{,}00518\cdot 10^{9}$ i~była identyczna dla każdej liczby lokacji od~1 do~9
(po dziesięć powtórzeń na~konfigurację), co empirycznie potwierdza niezależność
wyniku od~dekompozycji dziedziny.
Wymiana warstw brzegowych
jest więc niezbędna do~poprawności i~nie da się jej pominąć dla zysku
wydajności. Suma jest przy tym
miarą odporną na~błędy systematyczne, do~jakich należy pominięcie wymiany halo,
które rozchodzi pole w~jednym kierunku, a~nie na~błędy kompensujące się.
Zasada maksimum ogranicza każdą komórkę do~przedziału $[1,2]$, więc punktowe
naruszenie tej granicy zostałoby wychwycone niezależnie od sumy. Ta obowiązkowa
komunikacja, jak pokazano w~\cref{sec:skalowanie-komunikacja}, jest głównym wąskim
gardłem systemu.


\subsection{Skalowanie wątkowe na pojedynczym węźle}
\label{subsec:skalowanie-watkowe}

W~tej części przeanalizowano wydajność jednowęzłową. Celem tej serii eksperymentów
było zbadanie, jak czas wykonania programu zależy od liczby wątków obliczeniowych
przy ustalonym rozmiarze zadania, a~zatem ile faktycznie zyskuje się na
zrównolegleniu wewnątrzwęzłowym. Pomiary wykonano dla dziedziny $1000^3$ na~węźle klastra Pionier.

\Cref{tab:threads-1000} przedstawia wyniki tego eksperymentu. Faktyczna liczba
zadań qthreads jest tu równa liczbie żądanej.

\begin{table}[htbp]
\centering
\caption{Skalowanie wątkowe na pojedynczym węźle, dziedzina $1000^3$,
\texttt{-nl 1}, \num{10} kroków $\times$ \num{10} powtórzeń, Chapel 2.9,
kanał MPI, węzeł Pionier.}
\label{tab:threads-1000}
\begin{tabular}{rrrrrrrrr}
\toprule
Żądane & \texttt{maxTaskPar} & Mediana & Średnia & Min. & Maks. & Odch. std. & $S$ & $E$ \\
wątki & & [s] & [s] & [s] & [s] & [s] & & \\
\midrule
\num{1}  & \num{1}  & \num{9.812} & \num{9.814} & \num{9.803} & \num{9.829} & \num{0.008} & \num{1.000} & \num{1.000} \\
\num{2}  & \num{2}  & \num{9.948} & \num{10.104} & \num{9.941} & \num{10.470} & \num{0.240} & \num{0.986} & \num{0.493} \\
\num{4}  & \num{4}  & \num{8.362} & \num{8.362} & \num{8.356} & \num{8.365} & \num{0.003} & \num{1.173} & \num{0.293} \\
\num{8}  & \num{8}  & \num{8.094} & \num{8.095} & \num{8.087} & \num{8.101} & \num{0.004} & \num{1.212} & \num{0.152} \\
\num{16} & \num{16} & \num{8.739} & \num{8.741} & \num{8.731} & \num{8.763} & \num{0.008} & \num{1.123} & \num{0.070} \\
\bottomrule
\end{tabular}
\end{table}

Obraz jest tu jakościowo odmienny i~znacznie korzystniejszy. Przejście z~jednego do
dwóch wątków nie przynosi niemal zysku: mediana rośnie z~\SI{9.812}{\second} do
\SI{9.948}{\second}, a~trzy z~dziesięciu przebiegów dwuwątkowych były wyraźnie
wolniejsze, około \SI{10.47}{\second}, wskutek przejściowego ograniczenia taktowania.
Dopiero cztery wątki dają wyraźny skok wydajności: mediana spada do
\SI{8.362}{\second} (przyspieszenie \num{1.173}), a~optimum osiągane jest przy
ośmiu wątkach (\SI{8.094}{\second}, przyspieszenie \num{1.212}). Dalsze zwiększanie
liczby wątków ponownie szkodzi: przy szesnastu czas rośnie do
\SI{8.739}{\second}. Maksymalne uzyskane przyspieszenie wynosi zatem tylko około
\num{1.21}$\times$, mimo dostępności dwudziestu wątków.

Wydajność $E$ spada przy tym bardzo ostro: z~\num{1.000}
przez \num{0.493} i~\num{0.293} aż do \num{0.070} przy szesnastu wątkach. Tak
gwałtowny spadek to typowy objaw problemu
\emph{ograniczonego pasmem pamięci}~\cite{datta-stencil}. Schemat różnicowy jest operacją o~niskiej intensywności arytmetycznej, więc jego
przepustowość limituje pasmo pamięci głównej, a~nie liczba dostępnych rdzeni.
Po~nasyceniu tego pasma dodatkowe wątki rywalizują o~tę samą magistralę, co~prowadzi
do~degradacji wydajności.

Zwiększanie liczby wątków ponad pewien próg nie pomaga tu, a~wręcz szkodzi: po~nasyceniu
magistrali pamięci dodatkowe wątki nie znajdują użytecznej pracy.
\Cref{fig:threads-1000} ilustruje te zależności: lewy panel
przedstawia rozkład czasów wykonania w~postaci wykresu pudełkowego, a~prawy
zestawia zmierzone przyspieszenie z~przyspieszeniem idealnym oraz z~krzywą
wydajności, uwidaczniając odejście od skalowania liniowego już przy dwóch
wątkach.

\begin{figure}[htbp]
\centering
\includegraphics[width=0.95\linewidth]{threads_1000.pdf}
\caption{Skalowanie wątkowe dla dziedziny $1000^3$ na węźle Pionier.
Lewy panel: rozkład czasu wykonania w~postaci wykresu pudełkowego w~funkcji liczby
wątków. Prawy panel: zmierzone przyspieszenie na tle przyspieszenia idealnego
oraz wydajność $E$.}
\label{fig:threads-1000}
\end{figure}

W~części jednowęzłowej równoległość wątkowa przynosi w~najlepszym
razie skromny, około dwudziestoprocentowy zysk, a~wydajność szybko załamuje się
wraz ze wzrostem liczby wątków. Bezpośrednią przyczyną jest ograniczenie
pasmem pamięci. Sytuacja ta pogłębia się dopiero przy skalowaniu na~wiele węzłów.

\section{Skalowanie i ograniczenia komunikacji}
\label{sec:skalowanie-komunikacja}

W~tej części przeanalizowano zachowanie programu przy rozdzieleniu obliczeń na wiele
węzłów klastra połączonych siecią Ethernet \SI{1}{\giga\bit\per\second}.
Eksperymenty w~tej sekcji przeprowadzono na klastrze dla dziedziny $1000^3$
i~\num{100} kroków czasowych. Wykazano, że badany system jest \emph{ograniczony
komunikacją}: to przepustowość sieci, a~nie moc obliczeniowa węzłów, wyznacza
pułap osiąganej wydajności.

\subsection{Skalowanie liczbą węzłów}
\label{subsec:skalowanie-wezly}

\Cref{tab:node-scaling} przedstawia czas wykonania w~funkcji liczby węzłów, od
jednego do dziewięciu, dla dziedziny $1000^3$. Oprócz statystyk czasu podano
przyspieszenie, wydajność oraz objętość danych wymienianych w~ramach warstw
brzegowych na jeden krok czasowy w~kolumnie ,,Bajty/krok''.

\begin{table}[htbp]
\centering
\caption{Skalowanie liczbą węzłów, dziedzina $1000^3$, \num{100} kroków.}
\label{tab:node-scaling}
\begin{tabular}{lrrrrrrr}
\toprule
Węzły & Mediana & Średnia & Min. & Odch. std. & $S$ & $E$ & Bajty/krok \\
& [s] & [s] & [s] & [s] & & & \\
\midrule
\num{1}  & \num{87.344} & \num{87.307} & \num{87.104} & \num{0.092} & \num{1.000} & \num{1.000} & \num{0} \\
\num{2}  & \num{51.943} & \num{51.957} & \num{51.858} & \num{0.107} & \num{1.682} & \num{0.841} & \num{16000000} \\
\num{3}  & \num{52.978} & \num{53.413} & \num{52.862} & \num{1.316} & \num{1.649} & \num{0.550} & \num{32000000} \\
\num{4}  & \num{36.670} & \num{36.889} & \num{36.203} & \num{0.894} & \num{2.382} & \num{0.595} & \num{32032000} \\
\num{5}  & \num{41.476} & \num{41.482} & \num{41.347} & \num{0.093} & \num{2.106} & \num{0.421} & \num{64000000} \\
\num{6}  & \num{32.240} & \num{32.243} & \num{32.076} & \num{0.115} & \num{2.709} & \num{0.452} & \num{48064000} \\
\num{7}  & \num{37.114} & \num{37.112} & \num{36.909} & \num{0.156} & \num{2.353} & \num{0.336} & \num{96000000} \\
\num{8}  & \num{42.941} & \num{42.909} & \num{42.419} & \num{0.198} & \num{2.034} & \num{0.254} & \num{48096064} \\
\num{9}  & \num{36.832} & \num{36.817} & \num{36.469} & \num{0.157} & \num{2.371} & \num{0.263} & \num{64128000} \\
\bottomrule
\end{tabular}
\end{table}

Skalowanie węzłowe okazuje się wyraźnie \emph{niemonotoniczne}. Zamiast
oczekiwanego, systematycznego spadku czasu wraz ze wzrostem liczby węzłów,
obserwuje się silny ,,ząbek'' pomiędzy parzystymi a~nieparzystymi liczbami
węzłów. Najlepszy wynik uzyskano dla sześciu węzłów (mediana
\SI{32.240}{\second}, przyspieszenie \num{2.709}), podczas gdy sąsiednie
konfiguracje pięcio- i~siedmiowęzłowa wypadają zauważalnie gorzej
(\SI{41.476}{\second} oraz \SI{37.114}{\second}). Źródłem tego zjawiska jest
nierównomierny podział dziedziny $1000^3$: gdy liczba węzłów nie dzieli
gładko boku kostki, powstaje niezrównoważenie obciążenia, w~którym niektóre węzły
otrzymują większe poddziedziny i~stają się elementem spowalniającym całość,
co dodatkowo splata się ze wzrostem objętości komunikacji między węzłami.

Niezależnie od tych oscylacji, obraz ogólny jest jednoznaczny:
przyspieszenie \emph{płaszczy się} na poziomie około \num{2.4}$\times$ i~nie
zbliża się do skalowania liniowego, mimo dziewięciokrotnego zwiększenia liczby
węzłów~\cite{amdahl}. Wydajność spada z~\num{1.000} do zaledwie \num{0.263}
przy dziewięciu węzłach. Kluczem do wyjaśnienia jest ostatnia kolumna tabeli:
objętość wymienianych warstw brzegowych rośnie w~przybliżeniu liniowo wraz
z~liczbą węzłów, od \num{0} bajtów dla pojedynczego węzła, w~którym nie ma komunikacji
międzywęzłowej, do ponad \num{64000000} bajtów na krok przy dziewięciu węzłach.
Każdy dodany węzeł wnosi pewną moc obliczeniową, lecz zarazem powiększa
całkowity ruch sieciowy, a~przy powolnym łączu ten drugi efekt szybko
pochłania korzyść z~pierwszego. System jest zatem ograniczony komunikacją.

\begin{figure}[htbp]
\centering
\includegraphics[width=0.95\linewidth]{node_scaling.pdf}
\caption{Skalowanie liczbą węzłów dla dziedziny $1000^3$. Lewy panel: rozkład
czasu wykonania w~funkcji liczby węzłów. Prawy panel: przyspieszenie na tle
przyspieszenia idealnego oraz wydajność $E$.}
\label{fig:node-scaling}
\end{figure}

\begin{figure}[htbp]
\centering
\includegraphics[width=0.95\linewidth]{bytes_vs_nodes.pdf}
\caption{Objętość warstw brzegowych wymienianych na krok w~megabajtach oraz
mediana czasu wykonania w~funkcji liczby węzłów.}
\label{fig:bytes-vs-nodes}
\end{figure}

Zależności te zilustrowano na \cref{fig:node-scaling} oraz \cref{fig:bytes-vs-nodes}.
Pierwszy z~nich uwidacznia zarówno oscylacyjny charakter czasu wykonania, jak
i~znaczne odejście krzywej przyspieszenia od ideału liniowego. Drugi zestawia
rosnącą objętość komunikacji z~czasem wykonania, czyniąc bezpośrednio widocznym
związek między nasilającym się ruchem sieciowym a~stagnacją wydajności.

\subsection{Skalowanie wątkowe przy dziewięciu węzłach}
\label{subsec:skalowanie-watki-9n}

Aby rozstrzygnąć, czy wąskim gardłem w~konfiguracji wielowęzłowej jest
komunikacja, czy może niedostateczne wykorzystanie rdzeni w~obrębie węzła,
powtórzono skalowanie wątkowe przy ustalonej liczbie dziewięciu węzłów.
Wyniki zawiera \cref{tab:threads-9n}.

\begin{table}[htbp]
\centering
\caption{Skalowanie wątkowe przy dziewięciu węzłach, dziedzina $1000^3$,
\num{100} kroków.}
\label{tab:threads-9n}
\begin{tabular}{lrrrrrr}
\toprule
Konfiguracja & Mediana & Średnia & Min. & Odch. std. & $S$ & Obl./krok \\
& [s] & [s] & [s] & [s] & & [s] \\
\midrule
1 wątek   & \num{34.003} & \num{34.245} & \num{33.805} & \num{0.740} & \num{1.000} & \num{0.1422} \\
2 wątki   & \num{35.217} & \num{35.252} & \num{33.715} & \num{1.184} & \num{0.966} & \num{0.1481} \\
4 wątki   & \num{32.563} & \num{32.850} & \num{32.090} & \num{1.086} & \num{1.044} & \num{0.1600} \\
8 wątków  & \num{34.694} & \num{34.616} & \num{34.215} & \num{0.245} & \num{0.980} & \num{0.1775} \\
16 wątków & \num{36.623} & \num{36.633} & \num{36.208} & \num{0.195} & \num{0.928} & \num{0.1982} \\
\bottomrule
\end{tabular}
\end{table}

Odpowiedź jest jednoznaczna. Czas wykonania pozostaje niemal płaski w~zakresie
od około \num{32} do \num{37} sekund niezależnie od liczby wątków, z~płytkim
minimum przy czterech wątkach (\SI{32.563}{\second}). Przyspieszenie oscyluje
wokół jedności i~nigdy nie przekracza \num{1.044}. Najbardziej wymowna jest
jednak kolumna średniego czasu obliczeń na krok: rośnie ona
\emph{monotonicznie} wraz z~liczbą wątków, od \SI{0.1422}{\second} przy jednym
wątku do \SI{0.1982}{\second} przy szesnastu. Gdyby problem był ograniczony mocą
obliczeniową, dodawanie wątków skracałoby czas obliczeń na krok. Obserwuje się
jednak zjawisko przeciwne: czas ten rośnie. Nie jest to skutek nadmiernej liczby
wątków względem rdzeni, bo~te pozostają wolne ($p\le 16\le 20$), lecz rywalizacji
o~zasoby współdzielone: dodatkowe wątki nie znajdują użytecznej pracy, lecz
konkurują o~pasmo pamięci oraz z~wątkiem odpytującym warstwę komunikacyjną
GASNet~\cite{gasnet-spec}. Zwiększanie liczby wątków nie pomaga, ponieważ problem
jest ograniczony komunikacją, a~nie obliczeniami. \Cref{fig:threads-9n}
przedstawia te zależności graficznie.

\begin{figure}[htbp]
\centering
\includegraphics[width=0.95\linewidth]{threads_9nodes.pdf}
\caption{Skalowanie wątkowe przy dziewięciu węzłach, dziedzina $1000^3$.
Lewy panel: rozkład czasu wykonania. Prawy panel: średni czas obliczeń na krok
w~funkcji liczby wątków, rosnący wskutek rywalizacji o~zasoby współdzielone.}
\label{fig:threads-9n}
\end{figure}

\subsection{Skalowanie rozmiaru zadania}
\label{subsec:skalowanie-rozmiar}

Ostatni zestaw eksperymentów bada wpływ rozmiaru dziedziny na wydajność przy
ustalonej liczbie dziewięciu węzłów. \Cref{tab:cube-size} zestawia czasy
wykonania, przepustowość wyrażoną w~miliardach komórek na sekundę
(mld komórek/s) oraz udział komunikacji \texttt{comm\_frac}~\eqref{eq:commfrac}.

\begin{table}[htbp]
\centering
\caption{Skalowanie rozmiaru dziedziny przy dziewięciu węzłach, \num{100}
kroków.}
\label{tab:cube-size}
\begin{tabular}{lrrrr}
\toprule
Rozmiar & Mediana & Średnia & Przepustowość & \texttt{comm\_frac} \\
& [s] & [s] & [mld komórek/s] & \\
\midrule
$125^3$  & \num{1.668}   & \num{1.670}   & \num{0.117} & \num{0.963} \\
$250^3$  & \num{2.640}   & \num{2.652}   & \num{0.592} & \num{0.947} \\
$500^3$  & \num{7.625}   & \num{7.673}   & \num{1.639} & \num{0.729} \\
$1000^3$ & \num{36.741}  & \num{36.802}  & \num{2.722} & \num{0.454} \\
\bottomrule
\end{tabular}
\end{table}

Wynik ten najlepiej potwierdza, że ograniczenie ma charakter komunikacyjny.
Przepustowość rośnie systematycznie wraz z~rozmiarem dziedziny, od zaledwie
\SI{0.117}{} mld komórek/s dla najmniejszej kostki $125^3$ do \SI{2.722}{} mld komórek/s
dla największej kostki $1000^3$. Jednocześnie udział komunikacji
\texttt{comm\_frac} maleje monotonicznie z~\num{0.963} do \num{0.454}. Innymi
słowy, dla małych dziedzin blisko \SI{96}{\percent} czasu pochłania
komunikacja, a~faktyczne obliczenia stanowią margines. Wraz ze~wzrostem dziedziny
proporcje te wyrównują się i~przy~$1000^3$ komunikacja oraz obliczenia zajmują już
zbliżony czas. Nawet przy największej zbadanej kostce
komunikacja wciąż odpowiada za~blisko połowę czasu kroku. Trend jest wprawdzie
korzystny, lecz w~osiągalnym zakresie rozmiarów zadanie nie przestaje być istotnie
obciążone komunikacją.

Mechanizm tego zjawiska jest geometryczny: objętość dziedziny rośnie jak~$N^3$,
a~powierzchnia warstw brzegowych
jak~$N^2$, więc wraz ze wzrostem $N$ poprawia się stosunek pracy obliczeniowej
do~komunikacji. Duże kostki znacznie lepiej wykorzystują klaster niż małe,
zdominowane przez narzut komunikacyjny.

Rozmiar wymiany halo zależy od~geometrycznego układu dekompozycji. Dla
zilustrowania rozważmy dwuwymiarową dziedzinę $100\times 100$ podzieloną na cztery
lokale. Przy podziale plastrami wzdłuż jednej osi każdy blok ma wymiary
$25\times 100$, dwa skrajne lokale wysyłają po $100$ komórek brzegowych, a~dwa
wewnętrzne po $200$, co daje łącznie $600$ komórek wymienianych na krok. Przy
podziale na bloki kwadratowe $50\times 50$ każdy z~czterech lokali wysyła do
dwóch sąsiadów brzeg poziomy i~pionowy po $50$ komórek, co daje
$4\cdot(50+50)=400$ komórek. W~obu przypadkach objętość wymiany jest rzędu pola
powierzchni granicy, a~nie objętości, co w~trzech wymiarach przekłada się na
omówioną zależność $N^2$.

Największą zbadaną dziedziną jest $1000^3$. Omówione zależności ilustruje \cref{fig:cube-size}.

\begin{figure}[htbp]
\centering
\includegraphics[width=0.95\linewidth]{cube_size.pdf}
\caption{Skalowanie rozmiaru dziedziny przy dziewięciu węzłach. Lewy panel:
rozkład czasu wykonania dla kolejnych rozmiarów. Prawy panel: przepustowość
oraz udział komunikacji \texttt{comm\_frac} w~funkcji rozmiaru kostki.}
\label{fig:cube-size}
\end{figure}

\subsection{Synteza: system ograniczony komunikacją}
\label{subsec:synteza-comm}

Wyniki z tego rozdziału wskazują wspólnie, że system jest ograniczony
komunikacją i~przepustowością pamięci, a~nie mocą obliczeniową. Przyspieszenie
węzłowe płaszczy się na poziomie około \num{2.4}$\times$ (\cref{tab:node-scaling}),
daleko od skalowania liniowego, mimo dziewięciokrotnego zwiększenia liczby
węzłów. Zwiększanie liczby wątków na~węzeł nie przynosi poprawy, a~w~konfiguracji
wielowęzłowej średni czas obliczeń na~krok wręcz rośnie z~liczbą wątków
(\cref{tab:threads-9n}), co wskazuje na~nasycenie zasobów współdzielonych, a~nie na~niedobór
mocy obliczeniowej. Udział komunikacji dominuje przy małych dziedzinach i~maleje
dopiero dla dużych kostek (\cref{tab:cube-size}), zgodnie z~geometryczną relacją
między objętością $N^3$ a~powierzchnią $N^2$. Na każdym węźle jeden wątek warstwy
komunikacyjnej GASNet obciąża rdzeń w~niemal stu procentach, prowadząc aktywne
odpytywanie (ang.~\emph{busy-polling}) w~oczekiwaniu na~postęp komunikacji, co potwierdzono
analizą stosów wywołań w~debugerze. To nie jest obciążenie obliczeniowe, lecz
widoczny i~mierzalny skutek ograniczenia komunikacją.

Fizycznym ograniczeniem jest łącze Ethernet o~przepustowości
\SI{1}{\giga\bit\per\second}, które dla dziedziny rzędu $1000^3$ ulega
nasyceniu. Objętość warstw brzegowych na krok rośnie zarówno z~liczbą węzłów,
jak i~z~kwadratem boku dziedziny, wobec czego to sieć, a~nie procesory,
wyznacza pułap wydajności.

Osobnego omówienia wymaga dobór kanału komunikacyjnego GASNet, który okazał
się decydujący dla tego, czy obliczenia w~ogóle kończą się poprawnie~\cite{gasnet}.
W~pomiarach kanał \texttt{udp} przerywał działanie błędem \texttt{ECONGESTION}
pod obciążeniem typu incast, natomiast kanał \texttt{mpi} kończył obliczenia
poprawnie, choć wolniej. Z~tego powodu nowsza seria pomiarów wykorzystuje kanał
MPI: odporność na~zawodną sieć jest tu ważniejsza od szybkości.

Skoro system jest
ograniczony komunikacją, dalszej poprawy wydajności nie należy szukać ani
w~zwiększaniu liczby wątków, ani w~dokładaniu kolejnych węzłów, lecz przede
wszystkim w~przyspieszeniu warstwy komunikacyjnej, czy to przez zastosowanie
szybszego łącza sieciowego, czy przez wybór wydajniejszego kanału komunikacyjnego GASNet.
Dopóki halo wymieniane jest przez łącze \SI{1}{\giga\bit\per\second}, klaster
wciąż czeka na~dane z~sieci, a~część rdzeni pozostaje niewykorzystana.

\subsection{Zestawienie zbiorcze i~powtarzalność}
\label{subsec:synteza-tab}

Dla wygody czytelnika zebrano najważniejsze wnioski ilościowe z~poszczególnych
grup eksperymentów w~jednym miejscu. \Cref{tab:synteza} podaje dla każdej
serii kluczową miarę oraz jej najistotniejszy wynik.

\begin{table}[htbp]
\centering
\caption{Zestawienie zbiorcze najważniejszych wyników eksperymentalnych
rozdziału.}
\label{tab:synteza}
\begin{tabular}{@{}>{\raggedright\arraybackslash}p{4.2cm}>{\raggedright\arraybackslash}p{4.8cm}>{\raggedright\arraybackslash}p{4.2cm}@{}}
\toprule
Grupa eksperymentów & Kluczowa miara & Wynik \\
\midrule
Wątki $1000^3$ (\cref{tab:threads-1000})    & optimum / maks. przyspieszenie & 8 wątków, \num{1.21}$\times$ \\
Węzły $1000^3$ (\cref{tab:node-scaling})    & optimum / plateau przyspieszenia & 6 węzłów, około \num{2.4}$\times$ \\
Wątki, 9 węzłów (\cref{tab:threads-9n})     & przyspieszenie (płaskie)       & $\le \num{1.044}$ \\
Rozmiar kostki (\cref{tab:cube-size})       & przepustowość                  & \numrange{0.117}{2.722}~mld komórek/s \\
Rozmiar kostki (\cref{tab:cube-size})       & udział komunikacji \texttt{comm\_frac} & \num{0.963}$\to$\num{0.454} \\
\bottomrule
\end{tabular}
\end{table}

Z zestawienia widać, że niezależnie od osi skalowania, liczby wątków, liczby
węzłów czy rozmiaru zadania, wyniki trafiają na~to samo ograniczenie: pasmo
pamięci i~komunikację. Jedyną osią, wzdłuż której obserwuje się wyraźną poprawę
efektywności wykorzystania klastra, jest rozmiar zadania: powiększanie kostki
przesuwa proporcje na korzyść obliczeń, zgodnie z~omówioną wcześniej relacją
$N^3$ do $N^2$.

Odrębnego komentarza wymaga powtarzalność pomiarów, stanowiąca podstawę
zaufania do prezentowanych median. Zdecydowana większość serii charakteryzuje
się bardzo niskim rozrzutem wyników. Dla serii $1000^3$ na węźle
Pionier (\cref{tab:threads-1000}) odchylenie standardowe czasu wykonania
nie przekracza \SI{0.07}{\second}, a~dla większości konfiguracji jest rzędu
zaledwie kilku tysięcznych sekundy. Tak niski rozrzut oznacza, że pojedynczy przebieg jest
w~tych przypadkach niemal równie miarodajny co mediana z~dziesięciu powtórzeń,
a~różnice między konfiguracjami to rzędy dziesiątych sekundy lub więcej, więc nie da się ich wytłumaczyć samym szumem pomiarowym.

Osobną kwestią jest praktyczna granica rozmiaru zadania.
Dziedzin istotnie większych niż $1000^3$ nie udało się ukończyć: przekraczają
one możliwości pamięciowe i~czasowe klastra, wobec czego nie figurują w~żadnej
tabeli i~wyznaczają jedynie górny kres rozmiaru osiągalnego na badanej
infrastrukturze. Świadomość tego ograniczenia nie podważa głównych wniosków
rozdziału, które opierają się na~seriach o~niskim rozrzucie i~dziesięciokrotnej
powtarzalności, lecz trzeba o~niej pamiętać, oceniając, do~jakich rozmiarów
zadania sięgają wnioski tej pracy.
